\subsection{HD-TW01: averaging all tensor cycles without a field restriction}
\paragraph{Status.}
Complete root and separate /root/g34\_eightfolds whole-proof reviews
have passed, first in full message form during the outage and then on the
complete frozen TeX file after recovery. Source and exact finite-check
scopes are recorded separately; the final published4000 sketch is not
claimed as an independently recomputed certificate.
\paragraph{Actual-source hypotheses.}
Let an actual connected sphere source have a genuine generating product-one
inertia tuple \(T_1,\ldots,T_r\), period differential of rank nine, coefficient
degree \(e=a^t\), and transitive ordinary alternating core \(A_n^t\), with
\(t\ge2\) and \(a\ge4\).
Its core module is \(W^{\otimes t}\), where the compatible factor normalizers
act, up to scalars, through an irreducible \(S_n\)-module \(W\) of dimension \(a\).
Both the factor action on \(t\) letters and the natural block action on \(nt\)
letters are the actual quotient actions of this source.  In particular
\[
 r\ge12,\qquad D:=\sum_i\operatorname{rank}(1-T_i)=2e+6.
 \tag{TW.1}
\]
For a finite-order operator \(U\), write
\(\nu(U)=\dim U-\max_z\dim\ker(U-z)\).
Suppose that, for every nonidentity \(x\in S_n\),
\[
 \nu(W(x))\ge A\,\operatorname{ind}(x)+B
 \tag{TW.2}
\]
and that
\[
 \begin{gathered}
 A\ge0,\quad B>0,\quad d_0=A+B\ge1,\quad 2d_0\le a-1,\\
 A(n-1)+B\le a,\quad
 \lambda=1-\frac{nA+1}{a}\ge0,\quad
 L=a-(n-1)A\ge1,\quad a(9B-L-1)>6.
 \end{gathered}\tag{TW.3}
\]
Then the source does not exist.  This statement imposes no condition on the
imaginary quadratic field or scalar conductor.

\paragraph{Proof.}
A factor-fixing noncentral entry has \(m\ge1\) active factor operators
\(W(x_j)\), up to scalars.  Tensor spectral convolution bounds its normalized
projective support below by the maximum of the active normalized supports.
Averaging those \(m\) bounds and using \(m\le t\) gives
\[
 \frac{\operatorname{rank}(1-T_i)}e
 \ge \frac{A I_i}{ta}+\frac Ba ,
 \qquad I_i=\sum_j\operatorname{ind}(x_j).
 \tag{TW.4}
\]
A genuine scalar entry has actual moving rank \(e\), and \(I_i=0\), so it too
satisfies (TW.4).  This uses actual moving rank for such a scalar; its
projective support is zero.  Every genuine factor-fixing entry also has
moving rank at least \(d_0e/a\).
The accepted weighted-cycle minimum in D-ATT08 gives every factor-moving
entry moving rank at least \((a-1)e/(2a)\), hence at least \(d_0e/a\).
If four entries move factors, choose four of them and eight other genuine
entries.  Equation (TW.1) would imply
\[
 D\ge \frac{4(a-1)e}{2a}+\frac{8d_0e}a
   =2e+(8d_0-2)e/a
   \ge2e+6e/a>2e+6.
 \tag{TW.5}
\]
Thus at most three entries move factors, and at least nine fix them.

For any entry, let \(\pi\) be its factor permutation.  Its cycles have lengths
\(\ell_j\), with represented \(S_n\)-cycle products \(x_j\).  The exact
short-word estimate of D-ATT08 gives each cycle operator normalized support
at least
\[
 s_j\ge1-\frac1{\ell_j}
       +\frac{\nu(W(x_j))}{\ell_j a}
       -\frac{R_{\ell_j}(a)}{\ell_j a^{\ell_j}}.
 \tag{TW.6}
\]
Here \(R_1=0,\ R_2=a\), and
\(R_\ell(a)\le2a^{\ell-2}\) for \(\ell\ge3\), in the accepted range \(a\ge4\).
Consequently \(R_\ell(a)/a^\ell\le1/a\) for every \(\ell\ge2\).
Convolution across cycle operators gives total normalized support at least
\(\max_j s_j\), hence at least \(\sum_j(\ell_j/t)s_j\).
There are at most \(\operatorname{ind}(\pi)\) nontrivial cycles.
Also (TW.2) implies \(\nu(W(x))\ge A\operatorname{ind}(x)\) including \(x=1\).
The actual natural block index is
\[
 I_i=n\operatorname{ind}(\pi)+\sum_j\operatorname{ind}(x_j).
 \tag{TW.7}
\]
Substitution yields
\[
 \frac{\operatorname{rank}(1-T_i)}e
 \ge\frac{A I_i}{ta}
       +\frac{\lambda\,\operatorname{ind}(\pi)}t .
 \tag{TW.8}
\]
For a factor-fixing entry the additional \(B/a\) from (TW.4) remains available.
The two connected quotient covers satisfy
\(\sum_iI_i\ge2nt-2\) and
\(\sum_i\operatorname{ind}(\pi_i)\ge2t-2\).
Since \(A,\lambda\ge0\), summation gives, with \(r_{\rm fix}\ge9\),
\[
 \begin{split}
 D-2e
 &\ge\frac ea\left[r_{\rm fix}B-2-\frac2t(L-1)\right]\\
 &\ge\frac ea(9B-L-1)
 \ge a(9B-L-1)>6.
 \end{split}\tag{TW.9}
\]
This contradicts (TW.1), proving the assertion. \(\square\)

\paragraph{Scope of the mechanism.}
The weights \(\ell_j/t\) preserve the complete natural block index while
charging at most \(1/a\) for each nontrivial cycle.  No signed character floor
or distinguished scalar phase is a premise.  The proof applies only to
actual quotient actions and genuine source entries; it makes no existence
claim from a numerical profile.

\subsection{HD-TW02: an absolute-character criterion for every tensor count}
\paragraph{Statement.}
Let \(W\) be an irreducible \(S_n\)-module of degree \(a\), and let \(f(y)\)
denote the number of fixed letters of \(y\in S_n\).
Suppose \(n\ge24,\ a\ge n^3\), and
\[
 \frac{|\chi_W(y)|}{a}\le
       \left(\frac{\max\{f(y),1\}}n\right)^3
       \qquad(y\in S_n).
 \tag{TW.10}
\]
Then, for every nonidentity \(x\in S_n\),
\[
 \nu(W(x))\ge A\operatorname{ind}(x)+B,\qquad
 A=\frac{a(n-4)}{n^2},\quad B=\frac an.
 \tag{TW.11}
\]
No actual rank-nine sphere source with the compatible transitive core module
\(W^{\otimes t}\) exists, for any \(t\ge1\), over any imaginary quadratic
field, under the actual quotient hypotheses of HD-TW01.

\paragraph{Proof of the local bound.}
Put \(I=\operatorname{ind}(x)\), \(q=\operatorname{ord}(x)\),
and \(f_j=f(x^j)\), for \(0\le j<q\).
Fourier inversion on the cyclic group, followed by the triangle inequality
and (TW.10), gives
\[
 \frac{\nu(W(x))}{a}\ge
 1-\frac1q\sum_{j=0}^{q-1}
       \left(\frac{\max\{f_j,1\}}n\right)^3,\qquad
 \frac1q\sum_jf_j=n-I .
 \tag{TW.12}
\]
The second identity counts the orbits of \(\langle x\rangle\) on the letters.
If \(I\ge n/3\), use
\((\max\{f,1\}/n)^3\le f/n+n^{-3}\) to obtain
\[
 \nu(W(x))/a-AI/a
 \ge4I/n^2-n^{-3}
 \ge4/(3n)-n^{-3}>1/n.
 \tag{TW.13}
\]
If \(I<n/3\), every power has at least \(n-2I>n/3\) fixed letters:
the moved support of \(x\) has size at most \(2I\).
Thus \(u_j=f_j/n>1/3\), and
\(1-u_j^3\ge(13/9)(1-u_j)\).
For \(I\ge3\) this gives
\[
 \nu(W(x))/a-AI/a
 \ge 4I/(9n)+4I/n^2>1/n.
 \tag{TW.14}
\]
Finally, if \(I=1\) or \(2\), then \(u_j\ge1-4/n\ge5/6\), so
\(1+u_j+u_j^2\ge91/36>5/2\).
The resulting margin is at least
\(3I/(2n)+4I/n^2>1/n\).
These three cases prove (TW.11) for every target eigenvalue.

\paragraph{Actual-source contradiction.}
For the single factor the actual \(n\)-letter quotient and \(r\ge12\) give
\[
 D-2a\ge2A(n-1)+12B-2a
       =\frac{a(2n+8)}{n^2}>6 .
 \tag{TW.15}
\]
Genuine scalar entries also satisfy the required local rank bound:
their actual rank is \(a\), their natural index is zero, and \(B\le a\).
For \(t\ge2\), all hypotheses of HD-TW01 follow from
\[
 \begin{gathered}
 d_0=\frac{a(2n-4)}{n^2},\quad
 A(n-1)+B=a(1-2/n)^2,\quad
 \lambda=4/n-1/a,\\
 L=\frac{a(5n-4)}{n^2},\qquad
 9B-L-1=\frac{4a(n+1)}{n^2}-1 .
 \end{gathered}\tag{TW.16}
\]
Indeed \(d_0\ge1,\ L\ge1,\ \lambda>0\) because \(a\ge n^3\);
\(2d_0\le a-1\) follows from
\(a(n^2-4n+8)/n^2\ge n(n^2-4n+8)>1\).
The final expression in (TW.16), multiplied by \(a\), is larger than \(6\).
This proves the tensor assertion. \(\square\)

\paragraph{Published effective input.}
For \(n\ge4000\), the additional condition \(a\ge n^{120}\) implies
(TW.10).
The input is M\"uller--Schlage-Puchta,
\emph{Character theory of symmetric groups, subgroup growth of Fuchsian
groups, and random walks}, Advances in Mathematics 213 (2007), 919--982,
Theorem 2, with the exact normalization also stated in M\"uller,
\emph{Character theory of symmetric groups, analysis of long relators,
and random walks}, S\'eminaire Lotharingien de Combinatoire 56 (2006),
Article B56a, Section 5.1 equation (24), Theorem 9, and Lemma 7 equation (30).
The original author version locates Theorem 2 on page 8 and its explicit
\(4000\) discussion on pages 16--17.  That discussion is a proof sketch;
we use the published theorem, not an independently recomputed version of
its final numerical sketch.\cite{HD:MSP2007,HD:MullerSLC2006}

\paragraph{Proof of the stated consequence.}
For a nonidentity conjugacy class \(c\), let \(t_c\) be the least positive
even number of independent uniform elements of \(c\) having no common
fixed letter with probability at least \(1-1/n\).
The published theorem gives statistical mixing time at most \(10t_c\).
Its squared relative \(L^2\) normalization and character orthogonality give,
at the even time \(10t_c\),
\[
 \sum_{\chi(1)>1}\chi(1)^2
       \left|\frac{\chi(c)}{\chi(1)}\right|^{20t_c}
       \le \frac2n<1 .
 \tag{TW.17}
\]
If the first mixing time is earlier, the same bound at \(10t_c\) follows
because every normalized character ratio has modulus at most one.
For the particular degree \(a\), (TW.17) implies
\(|\chi_W(c)|/a<a^{-1/(10t_c)}\).

If \(1\le f=f(c)<n\), the elementary union bound for a common fixed letter
is \(n(f/n)^u\) at time \(u\).  Hence the even rounding gives
\[
 t_c\le2\left\lceil\frac{\log n}{\log(n/f)}\right\rceil
 \le\frac{4\log n}{\log(n/f)} .
 \tag{TW.18}
\]
It follows that
\[
 |\chi_W(c)|/a
 \le a^{-\log(n/f)/(40\log n)}
 \le (f/n)^3 .
 \tag{TW.19}
\]
For \(f=0\), \(t_c=2\), so the ratio is less than
\(a^{-1/20}\le n^{-6}\le n^{-3}\).
The identity has ratio one.  This proves (TW.10), including fixed-point-free
powers and the identity power, and finishes the consequence. \(\square\)

\subsection{HD-TW03: one effective ordinary alternating cutoff}
\paragraph{Statement and exact main scope.}
Set
\[
 \begin{gathered}
 M_* = 1+241\cdot240!\cdot2^{239},\qquad
 F_* = 240M_*+2,\\
 \mathcal N=\bigl[100(1+M_*)F_*^{240}\bigr]^2 .
 \end{gathered}\tag{TW.20}
\]
This is an explicit positive integer.  In the main nonsplit sixfold gate,
over a non-Eisenstein original imaginary quadratic field and within
the primitive factorwise tensor-indecomposable transitive-core scope
of D-BI19, no actual transitive ordinary alternating factor \(A_n\) has \(n\ge\mathcal N\).
All tensor counts and compatible scalar phases are included.
The earlier treatment of the single standard module and of small
alternating degrees is retained with its exact scope.
In particular the previously unbounded growing-tail ordinary alternating
range becomes finite.  This is a necessary-source reduction, not a new
detector-cycle construction or an executed enumeration.

\paragraph{A coefficient bound uniform over small tails.}
Let \(\mu\vdash k\), \(k\ge3\), and write the stable character polynomial as
\(q_\mu=Cx_1^k+R_\mu\), with \(C=f^\mu/k!\).
For the coefficient \(\ell^1\)-norm, the elementary binomial polynomial
satisfies
\[
 \left\|\binom{x}{m}\right\|_1=1\qquad(m\ge0).
 \tag{TW.21}
\]
For \(m\ge1\), its coefficients have alternating signs and their absolute
sum is
\(\prod_{j=0}^{m-1}(1+j)/m!=1\); the case \(m=0\) is constant one.

The Garsia--Goupil identity used in D-AFT01, equation (2.8),\cite{HD:GarsiaGoupil2009} expands
\(s_{(n-k,\mu)}\) as
\(\sum_{r=0}^k(-1)^r h_{n-k+r}e_r^\perp s_\mu\).
Here \(e_r^\perp s_\mu\) is the sum of \(s_\nu\) over vertical strips
\(\mu/\nu\) of size \(r\).
For \(j=k-r\), the associated induced character polynomial is
\[
 \sum_{\substack{m_1,m_2,\ldots\ge0\\ \sum i m_i=j}}
  \chi^\nu(1^{m_1}2^{m_2}\cdots)
  \prod_i\binom{x_i}{m_i}.
 \tag{TW.22}
\]
This is the invariant-\(j\)-subset character formula: each fixed subset
is a union of cycles, and its coefficient is the trace of its induced
\(S_j\)-action.  Equation (TW.21) and
\(|\chi^\nu|\le f^\nu\) bound its coefficient norm by \(p(j)f^\nu\).

By the ordinary branching rule,
\(\sum_{\mu/\nu\ {\rm vertical},\,|\nu|=j}f^\nu\le f^\mu\):
these are the dimensions of the sign-isotypic \(S_r\)-multiplicity spaces
inside the restriction of the \(S_k\)-module of degree \(f^\mu\).
Consequently
\[
 \|q_\mu\|_1
 \le f^\mu\sum_{j=0}^kp(j)
 \le f^\mu(k+1)2^{k-1}.
 \tag{TW.23}
\]
The last inequality uses \(p(0)=1\) and the injection of partitions of a
positive integer \(j\) into its \(2^{j-1}\) compositions.
The AFT coefficient parameter therefore satisfies
\[
 M_\mu:=1+C^{-1}\|R_\mu\|_1
 \le 1+k!(k+1)2^{k-1}.
 \tag{TW.24}
\]
Thus, for \(3\le k\le240\),
\(M_\mu\le M_*\), \(C^{-1}\le240!\), and
\(F_\mu=\lceil kM_\mu\rceil+1\le F_*\).
Substitution into the full accepted threshold
\[
 N_\mu=\left\lceil\max\left\{
 2k+2,200,100/C,100kM_\mu,
 [100(1+M_\mu)]^5,
 [100(1+M_\mu)F_\mu^k]^2
 \right\}\right\rceil
 \tag{TW.25}
\]
shows \(N_\mu\le\mathcal N\).
For clarity, \(F_*\ge100(1+M_*)\), so the last term of (TW.20) dominates
also the fifth power and every displayed linear term.  It also exceeds
\(482^2=232324\) and \(4000\).

\paragraph{The fixed-tail estimates in every tensor count.}
For \(n\ge N_\mu\), D-AFT01--02 gives
\(S=Cn^{k-1}\ge100\),
\(.99nS\le a\le1.01nS\),
\[
 A=a/n-101S/(100n),\quad B=2S/5,\quad
 .9S\le A\le1.01S,\quad1.3S\le A+B\le1.41S.
 \tag{TW.26}
\]
These are precisely local all-eigenvalue bounds, so they are valid before
any conductor argument.
For HD-TW01,
\(\lambda=(1.01S-1)/a>0\),
\[
 L=a/n+1.01S(1-1/n),\qquad
 1\le L\le2.02S,\qquad
 9B-L-1\ge1.58S-1>0.
 \tag{TW.27}
\]
Since \(n\ge200\), one has
\(2(A+B)\le2.82S\le a-1\),
\(L\ge .99S>B\), and \(a(1.58S-1)>6\).
Thus all remaining hypotheses of HD-TW01 hold.
It follows that these same thresholds exclude every transitive tensor
count over every imaginary quadratic field.
The single-factor contradiction is already D-AFT02.
This extension does not change the separate standard-module premise.

\paragraph{The growing-tail range.}
Suppose now that an actual source in the stated main gate has
\(n\ge\mathcal N\).
The accepted HD-ALT-07/08 shape reduction, after transposing the partition
if necessary, gives an ordinary label \(\lambda=(n-k,\mu)\) with
\(k<n/4\).  The already accepted standard and two-box exclusions, together
with the three-box exclusions, leave \(k\ge4\).
If \(k\le240\), the preceding paragraph excludes the source.

If \(k\ge241\), the published degree bound of
M\"uller--Schlage-Puchta, Lemma 8(ii),\cite{HD:MSP2007} gives
\[
 a=f^\lambda\ge\binom{n-k}{k}f^\mu
              \ge\binom{n-k}{k}.
 \tag{TW.28}
\]
Let \(N_0=\lfloor3n/4\rfloor\).
Then \(n-k\ge N_0\), \(241\le k\le N_0/2\), and \(N_0\ge n/2\).
Monotonicity of binomial coefficients through the midpoint, followed by
the elementary product inequality
\(\binom{N_0}{241}\ge(N_0/241)^{241}\), gives
\[
 a\ge\binom{N_0}{241}
   \ge(n/482)^{241}
   \ge n^{120+1/2}\ge n^{120}.
 \tag{TW.29}
\]
Here \(n\ge482^2\) justifies the penultimate inequality.
The published effective mixing input and HD-TW02 exclude this case as well.
Every possible tail is covered, proving \(n<\mathcal N\). \(\square\)

\paragraph{Geometric and programme meaning.}
The indices in HD-TW01 are ramification indices of actual quotient covers
of the original Hurwitz source; the contradiction is a strengthened
Chevalley--Weil Euler budget for its weight-one isotypic local system.
The character-polynomial calculation uses induction and exterior-power
operations in the representation \(\lambda\)-ring.  It passes to the
corresponding group-algebra summands of curve cohomology and their
realizations, but asserts no additional algebraic cycle, motivic splitting,
weight argument, or special-value bound for a motivic \(L\)-function.
Together with HD-TC01--03, the now finite ordinary alternating catalogue
has a finite effective inducing-image and component-rank search.
Original-\(K\) descent and the actual rational polarization remain separate.

For Grothendieck's programme, this is a strict reduction of the possible
curve suppliers for the integral uniform form of condition A
(cote16, scan050, corrected in SIGN_NEXT SN-C03/C04).
The positivity of condition I (scan049 and cote44, Theorem4) still cannot
supply missing cycles.  The projector/division mechanism remains the
Cayley--Hamilton inverse in scans048--049 and the already accepted run
statements.  No use of Markman, Severn, or032 is a premise of HD-TW01--03.
A counterexample to the eventual uniform detector bound would still be
a ghost in the precise GHOSTS B03 sense.

